\documentclass[11pt]{article}
\usepackage[margin=1in]{geometry}
\usepackage{setspace}
\usepackage{caption}
\usepackage{booktabs}
\usepackage{graphicx}
\usepackage{hyperref}
\usepackage{amsmath}
\usepackage{threeparttable}

\onehalfspacing

% Customize section, figure and table numbering to S1, S2, ...
\renewcommand{\thesection}{S\arabic{section}}
\renewcommand{\thefigure}{S\arabic{figure}}
\renewcommand{\thetable}{S\arabic{table}}
\captionsetup[figure]{name={Figure}}
\captionsetup[table]{name={Table}}

% % Hyperlink setup
% \hypersetup{
%     colorlinks=true,
%     linkcolor=black,
%     filecolor=magenta,      
%     urlcolor=cyan,
% }

\usepackage[round]{natbib} 
\makeatletter
\renewcommand{\bibsection}{}
\makeatother

% Title with main paper as subtitle
\title{
  \textbf{Supplementary Materials for} \\
  \vspace{1em}
  Causal discovery in urban data with temporal empirical dynamic modeling: The R package tEDM
}
\author{Lv et al.}
\date{}

\begin{document}

\maketitle

% Table of contents
% \newpage
\tableofcontents


% List of figures
\newpage
\listoffigures
% \newpage

% List of tables
\listoftables

\newpage
\section{Benchmarking results}\label{sec_benchmarks}

In total, five benchmarks were simulated in this study. The initial spatial cross sections were generated by drawing multivariate normal samples using the \texttt{MASS} R package \citep{rpkg_MASS}, and the simulated outputs were stored as spatial raster data using the \texttt{terra} package \citep{rpkg_terra}. 

Section~\ref{sec_main_benchmarks} reports the detailed cross mapping results for three representative causal structures reflecting key inferential challenges in multivariate spatial cross sectional data: an indirect-only pathway, a configuration where direct and indirect causal influences jointly reinforce the same direction, and a configuration where the direct and indirect causal influences oppose each other. Section~\ref{sec_collider_benchmarks} presents results for the additional scenario involving a collider structure, including the simulated spatial cross sections and cross-prediction diagnostics. Section~\ref{sec_confounding_benchmarks} provides corresponding results for the confounding structure. Across all five benchmarks, the minimum embedding dimension was determined using the false nearest neighbor method \citep{kennel1992fnn}, and the optimal number of nearest neighbors in the simplex projection was selected based on self-mapping performance of each variable \citep{sugihara1990simplex_projection} .


\subsection{Performance across canonical causal structures}\label{sec_main_benchmarks}

\paragraph{Cross mapping for simulated benchmark: $X \rightarrow Y \rightarrow Z$}

Figure~\ref{fig_benchmark_s1} presents the detailed cross mapping results for benchmark scenario~1 where SCPCM correctly recovers the mediated causal chain $X \rightarrow Y \rightarrow Z$. Panels~a, c, and e of Figure~\ref{fig_benchmark_s1} show GCCM results, while panels~b, d, and f display SCPCM outputs. In Figure~\ref{fig_benchmark_s1}b, the conditioned cross mapping $Y \; xmap \; X \mid Z$ converges to $\rho = 0.36$ at the largest library size, indicating a direct causal influence from $X$ to $Y$. Conversely, $X \; xmap \; Y \mid Z$ exhibits only a shallow increase and remains weak at $\rho = 0.13$, suggesting no direct causal influence from $Y$ to $X$. By comparison, Figure~\ref{fig_benchmark_s1}a shows that GCCM correctly detects an $X \rightarrow Y$ influence but also produces a spurious inference of a weak reverse influence ($Y \rightarrow X$). For the $Y \text{-} Z$ pair, both GCCM (Figure~\ref{fig_benchmark_s1}c) and SCPCM (Figure~\ref{fig_benchmark_s1}d) identify that $Z \; xmap \; Y$ continues to increase with library size, whereas the reverse direction converges early, supporting a direct causal link $Y \rightarrow Z$. For the $X \text{-} Z$ pair, GCCM suggests bidirectional causality (Figure~\ref{fig_benchmark_s1}e). In contrast, SCPCM (Figure~\ref{fig_benchmark_s1}f) shows that conditioning on $Y$ collapses the predictive skill of $Z \; xmap \; X$ to the same low level as $X \; xmap \; Z$, while the trend of $X \; xmap \; Z$ remains stable before and after conditioning. This indicates the absence of any direct causal influence between $X$ and $Z$.

\begin{figure}[htbp]
\centering
\includegraphics[width=0.9\textwidth]{figure/SI_benchmark1.png} 
\caption{Cross mapping results for benchmark scenario~1.}
\label{fig_benchmark_s1}
\end{figure}

\newpage
\paragraph{Cross mapping for simulated benchmark: $X \rightarrow Y \rightarrow Z,\; X \rightarrow Z$} 

Figure~\ref{fig_benchmark_s2} presents the cross mapping results for benchmark scenario~2. In Figure~\ref{fig_benchmark_s2}b, the conditioned cross mapping $Y \; xmap \; X \mid Z$ continues to increase with increasing library size, indicating a direct causal influence from $X$ to $Y$. By contrast, $X \; xmap \; Y \mid Z$ converges early to a low level ($\rho = 0.28$), implying the absence of a direct causal influence from $Y$ to $X$. GCCM similarly detects an $X \rightarrow Y$ influence (Figure~\ref{fig_benchmark_s2}a), but again produces a spurious inference of a weak reverse influence. For the $Y \text{-} Z$ pair, SCPCM clearly identifies a direct causal link $Y \rightarrow Z$, while excluding direct causal influence of $Z$ on $Y$ (Figure~\ref{fig_benchmark_s2}d), as the $Z \; xmap \; Y \,|\, X$ curve remains negative and statistically nonsignificant. GCCM similarly detects an $Y \rightarrow Z$ causal influence (Figure~\ref{fig_benchmark_s2}c) (Figure~\ref{fig_benchmark_s2}c). For the $X \text{-} Z$ pair, GCCM suggests reciprocal causation (Figure~\ref{fig_benchmark_s2}e). However, SCPCM reveals a different conclusion. The conditioned predictability $Z \; xmap \; X \mid Y$ exhibits smooth convergence (final $\rho = 0.38$), whereas $X \; xmap \; Z \mid Y$ shows a shallow rise followed by a decline (Figure~\ref{fig_benchmark_s2}f). SCPCM therefore indicates that only a direct causal influence $X \rightarrow Z$ exists, while the reverse direction $Z \rightarrow X$ is not supported.

\begin{figure}[htbp]
\centering
\includegraphics[width=0.9\textwidth]{figure/SI_benchmark2.png}
\caption{Cross mapping results for benchmark scenario~2.}
\label{fig_benchmark_s2}
\end{figure}

\newpage
\paragraph{Cross mapping for simulated benchmark: $X \rightarrow Y \rightarrow Z \rightarrow X$}

Figure~\ref{fig_benchmark_s3} presents the cross mapping results for benchmark scenario~3. For the $X \text{-} Y$ pair, GCCM detects a dominant $X \rightarrow Y$ causal influence but simultaneously infers a spurious weak reverse causal influence (Figure~\ref{fig_benchmark_s3}a). In contrast, SCPCM shows that $X \; xmap \; Y \mid Z$ converges slowly to a low level ($\rho = 0.21$), whereas $Y \; xmap \; X \mid Z$ maintains a clear increasing trend and consistently exceeds the reverse direction (Figure~\ref{fig_benchmark_s3}b). These results indicate that only a direct causal influence from $X$ to $Y$ exists. For the $Y \text{-} Z$ pair, GCCM suggests a unidirectional influence from $Y$ to $Z$ (Figure~\ref{fig_benchmark_s3}c). SCPCM further clarifies the directional causation: the $Y \; xmap \; Z \mid X$ curve shows a shallow rise followed by a decline, while the $Z \; xmap \; Y \mid X$ curve continues to increase throughout the library expansion (Figure~\ref{fig_benchmark_s3}d). These results support a direct causal influence $Y \rightarrow Z$, with no direct causal influence in the reverse direction. For the $X \text{-} Z$ pair, GCCM again implies reciprocal causation (Figure~\ref{fig_benchmark_s3}e). SCPCM instead reveals strong and sustained growth in $Z \; xmap \; X \mid Y$, whereas $X \; xmap \; Z \mid Y$ shows only a weak rise followed by a pronounced decline (Figure~\ref{fig_benchmark_s3}f). Thus, SCPCM infers a direct causal influence $Z \rightarrow X$, while excluding any direct causal influence from $X$ to $Z$.

\begin{figure}[htbp]
\centering
\includegraphics[width=0.9\textwidth]{figure/SI_benchmark3.png}
\caption{Cross mapping results for benchmark scenario~3.}
\label{fig_benchmark_s3}
\end{figure}

\newpage
\subsection{Robustness to spurious associations induced by collider structures}\label{sec_collider_benchmarks}

We further evaluate the capability of SCPCM to suppress spurious causal inference in a collider structure $X \rightarrow Y \leftarrow Z$. The spatial distributions of $X$, $Y$, and $Z$ under this structure are shown in Figure~\ref{fig_benchmark_s4}a. When applying GCCM, both directions of cross prediction between $X$ and $Z$ indicate increasing predictability (Figure~\ref{fig_benchmark_s4}b), leading GCCM to incorrectly infer reciprocal causal influence between the two variables. In contrast, SCPCM reveals a distinctly different behavior: once conditioning on the collider variable $Y$, the cross-mapping skills of $X \; xmap \; Z \mid Y$ and $Z \; xmap \; X \mid Y$ become nearly identical as the library size approaches the full sample (Figure~\ref{fig_benchmark_s4}c). Together with the unconditioned GCCM results, this convergence to symmetry under conditioning provides strong evidence that no direct causal influence exists between $X$ and $Z$.

\begin{figure}[htbp]
\centering
\includegraphics[width=1\textwidth]{figure/SI_benchmark4.png}
\caption{Simulated spatial cross-section of a collider structure and its corresponding cross mapping results.}
\label{fig_benchmark_s4}
\end{figure}

\newpage
\subsection{Robustness to false positives under controlled confounding}\label{sec_confounding_benchmarks}

We also evaluate the capability of SCPCM to suppress spurious causal inference between $X$ and $Z$ in a confounding structure $X \leftarrow Y \rightarrow Z$. The spatial distributions of $X$, $Y$, and $Z$ in this setting are shown in Figure~\ref{fig_benchmark_s5}a. Under GCCM, both cross-prediction directions between $X$ and $Z$ exhibit strong and increasing predictability (Figure~\ref{fig_benchmark_s5}b), causing GCCM to incorrectly infer reciprocal causal influence between $X$ and $Z$. In contrast, when conditioning on the confounder $Y$, both $X \; xmap \; Z \mid Y$ and $Z \; xmap \; X \mid Y$ initially show shallow increases in predictability followed by declines as the library size grows (Figure~\ref{fig_benchmark_s5}c). Both cross prediction directions toward similarly low skill demonstrates that $X$ and $Z$ share predictive information only through their mutual dependence on $Y$, and therefore no direct causal influence exists between them.

\begin{figure}[htbp]
\centering
\includegraphics[width=1\textwidth]{figure/SI_benchmark5.png}
\caption{Simulated spatial cross-section of a confounding structure and its corresponding cross mapping results.}
\label{fig_benchmark_s5}
\end{figure}

\newpage
\section{Comparative performance on real-world geospatial datasets}
\label{sec_cases}

Two real world applications of SCPCM are presented in this study, alongside comparisons with Pearson correlation \citep{pearson1895}, partial correlation \citep{baba2004partialcor} , DirectLiNGAM \citep{shimizu2011directlingam} , and GCCM \citep{gao2023gccm} . SCPCM and GCCM analyses were conducted using the \texttt{spEDM} R package developed in this work \citep{rpkg_spEDM}. Pearson correlation coefficients were computed using the \texttt{psych} R package \citep{rpkg_psych}, partial correlations using the \texttt{ppcor} R package \citep{rpkg_ppcor}, and DirectLiNGAM results using the \texttt{lingam} python package \citep{pypkg_lingam}.

\subsection{Farmland net primary productivity case}\label{sec_case1}

The first application examines SCPCM in recovering the mediated causal pathway among elevation, temperature, and farmland NPP in mainland China. Temperature, precipitation, and NPP were obtained and preprocessed following the same data sources and procedures as in the original study \citep{gao2022ts_sp_causality,gao2023gccm}, while SRTM Digital Elevation Data Version~4 were extracted using the \texttt{rgee} package from the Google Earth Engine platform to provide elevation information \citep{rpkg_rgee}. All raster datasets were reprojected and harmonized to a common coordinate system using the \texttt{terra} package, and resampled to a spatial resolution of 30~km \citep{rpkg_terra}. The spatial distributions of farmland elevation, NPP, temperature, and precipitation in mainland China were visualized using the \texttt{ggplot2} and \texttt{tidyterra} R packages \citep{rpkg_ggplot2,rpkg_tidyterra}. For GCCM and SCPCM, 1,500 non-NaN grid cells were randomly selected as the prediction set, and the remaining valid cells were used as the library set for cross prediction. Figure~\ref{fig_cor_case1}a and Figure~\ref{fig_cor_case1}b present the Pearson and partial correlation results among the variables in this application, respectively. Figure~\ref{fig_cross_mapping_case1} supplements the results of GCCM by showing the cross prediction performance of precipitation and temperature on NPP. A complete summary of the results for this application is provided in Table~\ref{tab_case1}.

% \newpage
\subsection{Infrastructure accessibility and health outcome case}\label{sec_case2}

The second application investigates the direct causal influences of economic, social, and environmental infrastructure accessibility on health outcomes. The raw indicators were obtained from the original study \citep{tu2025infra_health} and spatially joined with national administrative boundaries from the \texttt{geocn} R package using the \texttt{sf} package \citep{rpkg_geocn,rpkg_sf1,rpkg_sf2}, resulting in a polygon‐based spatial vector dataset. Countries without any land neighbors were removed and spatial neighborhood structures were then constructed using the \texttt{spdep} package to support GCCM and SCPCM analyses \citep{rpkg_spdep}. The spatial distributions of all four variables involved in this application were visualized using the \texttt{tmap} package \citep{rpkg_tmap}. Figure~\ref{fig_cor_case2}a and Figure~\ref{fig_cor_case2}b present the Pearson and partial correlation results among the variables in this application, respectively. Figure~\ref{fig_cross_mapping_case2} further supplements the results of GCCM in this application. A complete summary of the results for this application is provided in Table~\ref{tab_case2}.

\begin{figure}[htbp]
\centering
\includegraphics[width=1\textwidth]{figure/SI_elev_npp_cor.png}
\caption{Pairwise pearson and partial correlations between elevation (elev), temperature (tem), precipitation (pre) and farmland NPP (npp).}
\label{fig_cor_case1}
\end{figure}

\begin{figure}[htbp]
\centering
\includegraphics[width=0.8\textwidth]{figure/SI_elev_npp.png}
\caption{Cross mapping results between farmland NPP and climatic drivers (temperature and precipitation) without conditioning variables.}
\label{fig_cross_mapping_case1}
\end{figure}

\begin{table}[htbp]
\centering
\caption{Comparison of causal discovery results for the farmland NPP case.}
\label{tab_case1}
\begin{threeparttable}
\begin{tabular}{@{}llcc@{}}
\toprule
Method & Relation & Strength & Significance \\
\midrule
Pearson Correlation & Elev $\leftrightarrow$ NPP & $-0.12$ & -- \\
 & Pre $\leftrightarrow$ NPP & $0.79$ & -- \\
 & Tem $\leftrightarrow$ NPP & $0.66$ & -- \\
 & Elev $\leftrightarrow$ Tem & $-0.34$ & -- \\
 & Elev $\leftrightarrow$ Pre & $-0.28$ & -- \\
 & Pre $\leftrightarrow$ Tem & $0.78$ & -- \\
\addlinespace

Partial Correlation & Elev $\leftrightarrow$ NPP & $0.20$ & -- \\
 & Pre $\leftrightarrow$ NPP & $0.60$ & -- \\
 & Tem $\leftrightarrow$ NPP & $0.16$ & -- \\
 & Elev $\leftrightarrow$ Tem & $-0.22$ & -- \\
 & Elev $\leftrightarrow$ Pre & $-0.14$ & -- \\
 & Pre $\leftrightarrow$ Tem & $0.50$ & -- \\
\addlinespace

DirectLiNGAM & Pre $\rightarrow$ NPP & $0.38$ & -- \\
 & NPP $\rightarrow$ Elev & $0.12$ & -- \\
 & Elev $\rightarrow$ Tem & $-0.01$ & -- \\
 & Pre $\rightarrow$ Tem & $0.01$ & -- \\
 & Pre $\rightarrow$ Elev & $-0.10$ & -- \\
\addlinespace

GCCM & Elev \; xmap \; NPP & $0.58$ & -- \\
 & NPP \; xmap \; Elev & $0.74$ & -- \\
 & Pre \; xmap \; NPP & $0.65$ & -- \\
 & NPP \; xmap \; Pre & $0.79$ & -- \\
 & Tem \; xmap \; NPP & $0.53$ & -- \\
 & NPP \; xmap \; Tem & $0.73$ & -- \\
\addlinespace

SCPCM & Elev \; xmap \; NPP $|$ Pre & $0.53$ & -- \\
 & NPP \; xmap \; Elev $|$ Pre & $0.71$ & -- \\
 & Elev \; xmap \; NPP $|$ Tem & $0.53$ & -- \\
 & NPP \; xmap \; Elev $|$ Tem & $0.47$ & -- \\
 & Elev \; xmap \; NPP $|$ Pre, Tem & $0.50$ & -- \\
 & NPP \; xmap \; Elev $|$ Pre, Tem & $0.45$ & -- \\
\bottomrule
\end{tabular}
\begin{tablenotes}
     \footnotesize
     \item Note: ``$\leftrightarrow$'' denotes a symmetric association estimated by Pearson or partial correlation; ``$\rightarrow$'' indicates a direct causal influence inferred by DirectLiNGAM (relations not reported by DirectLiNGAM are deemed non-causal under its model assumptions); Cross prediction is denoted as ``$xmap$'' . For SCPCM, conditioning variables are listed after the vertical bar. All $p$-values are $< 0.001$ (denoted by ``--'').
\end{tablenotes}
\end{threeparttable}
\end{table}

\begin{figure}[htbp]
\centering
\includegraphics[width=1\textwidth]{figure/SI_infra_health_cor.png}
\caption{Pairwise pearson and partial correlations between three types of infrastructure accessibility and health outcome.}
\label{fig_cor_case2}
\end{figure}

\begin{figure}[htbp]
\centering
\includegraphics[width=1\textwidth]{figure/SI_infra_health.png}
\caption{Cross mapping results between three types of infrastructure accessibility and health outcome without conditioning variables}
\label{fig_cross_mapping_case2}
\end{figure}

\begin{table}[htbp]
\centering
\caption{Comparison of causal discovery results for the infrastructure accessibility and health outcome case.}
\label{tab_case2}
\begin{threeparttable}
\begin{tabular}{@{}llcc@{}}
\toprule
Method & Relation & Strength & Significance \\
\midrule
Pearson Correlation & Eco $\leftrightarrow$ HALE & $0.52$ & -- \\
 & Soc $\leftrightarrow$ HALE & $0.48$ & -- \\
 & Env $\leftrightarrow$ HALE & $0.27$ & 0.003 \\
 & Eco $\leftrightarrow$ Soc & $0.48$ & -- \\
 & Eco $\leftrightarrow$ Env & $-0.12$ & 0.146 \\
 & Soc $\leftrightarrow$ Env & $0.20$ & 0.036 \\
\addlinespace

Partial Correlation & Eco $\leftrightarrow$ HALE & $0.44$ & -- \\
 & Soc $\leftrightarrow$ HALE & $0.22$ & 0.008 \\
 & Env $\leftrightarrow$ HALE & $0.33$ & -- \\
 & Eco $\leftrightarrow$ Soc & $0.36$ & -- \\
 & Eco $\leftrightarrow$ Env & $-0.36$ & -- \\
 & Soc $\leftrightarrow$ Env & $0.20$ & 0.018 \\
\addlinespace

DirectLiNGAM & Eco $\rightarrow$ HALE & $10.55$ & 0.038 \\
 & Soc $\rightarrow$ HALE & $4.73$ & 0.004 \\
 & Env $\rightarrow$ HALE & $6.52$ & 0.712 \\
 & Soc $\rightarrow$ Eco & $0.47$ & -- \\
 & Eco $\rightarrow$ Env & $-0.29$ & 0.006 \\
 & Soc $\rightarrow$ Env & $0.33$ & 0.348 \\
\addlinespace

GCCM & Eco \; xmap \; HALE & $0.28$ & -- \\
 & HALE \; xmap \; Eco & $0.38$ & -- \\
 & Soc \; xmap \; HALE & $0.44$ & -- \\
 & HALE \; xmap \; Soc & $0.18$ & 0.030 \\
 & Env \; xmap \; HALE & $0.38$ & -- \\
 & HALE \; xmap \; Env & $0.50$ & -- \\
\addlinespace

SCPCM & Eco \; xmap \; HALE $|$ Soc, Env & $0.01$ & 0.860 \\
 & HALE \; xmap \; Eco $|$ Soc, Env & $0.22$ & 0.007 \\
 & Soc \; xmap \; HALE $|$ Eco, Env & $0.09$ & 0.305 \\
 & HALE \; xmap \; Soc $|$ Eco, Env & $0.13$ & 0.123 \\
 & Env \; xmap \; HALE $|$ Eco, Soc & $0.07$ & 0.440 \\
 & HALE \; xmap \; Env $|$ Eco, Soc & $0.35$ & -- \\
\bottomrule
\end{tabular}
\begin{tablenotes}
     \footnotesize
     \item Note: ``$\leftrightarrow$'' denotes a symmetric association estimated by Pearson or partial correlation; ``$\rightarrow$'' indicates a direct causal influence inferred by DirectLiNGAM (relations not reported by DirectLiNGAM are deemed non-causal under its model assumptions); Cross prediction is denoted as ``$xmap$'' . For SCPCM, conditioning variables are listed after the vertical bar. $p$-values less than $0.001$ are denoted by ``--”.
\end{tablenotes}
\end{threeparttable}
\end{table}

\newpage
\section{Sensitivity analysis of SCPCM to observational noise}\label{sec_sensitivity}

% \subsection{Sensitivity to observational noise}
% \label{sec_noise}

We evaluated the ability of SCPCM to handle observational noise across the three scenarios described in Section~\ref{sec_main_benchmarks}. Specifically, we added random white noise after each interaction in Eq. 2 of the main text. The updated system is

\begin{equation}\label{add_white_noises}
\left\{
\begin{aligned}
x_{i}^{(t+1)} &= 1 - \alpha_x\, x_i^{(t)} \left( \frac{1}{k} \sum_{j \in \mathcal{N}(i)} x_j^{(t)} - \beta_{yx}\, y_i^{(t)} - \beta_{zx}\, z_i^{(t)} \right) + \xi_{x,i}^{(t)} \\
y_{i}^{(t+1)} &= 1 - \alpha_y\, y_i^{(t)} \left( \frac{1}{k} \sum_{j \in \mathcal{N}(i)} y_j^{(t)} - \beta_{xy}\, x_i^{(t)} - \beta_{zy}\, z_i^{(t)} \right) + \xi_{y,i}^{(t)} \\
z_{i}^{(t+1)} &= 1 - \alpha_z\, z_i^{(t)} \left( \frac{1}{k} \sum_{j \in \mathcal{N}(i)} z_j^{(t)} - \beta_{xz}\, x_i^{(t)} - \beta_{yz}\, y_i^{(t)} \right) + \xi_{z,i}^{(t)} 
\end{aligned}
\right.
\end{equation}

Here, $\xi_{x,i}^{(t)}$, $\xi_{y,i}^{(t)}$ and $\xi_{z,i}^{(t)}$ denote white noise with zero mean and varying standard deviations. These standard deviations represent different levels of observational noise, that is, different signal to noise ratios in this approach. Following earlier work \citep{leng2020pcm,tao2023cmc}, and given that our simulations use a different spatial logistic map \citep{willeboordse2003spatial_logistic_map}, we derived the baseline noise scale from the sample variance of the benchmark scenario without added white noise. We then scaled this baseline by 15\% , 30\% , 45\% , 60\% , 75\%  and 90\% to generate the standard deviations of the added white noise. SCPCM was run with the same parameter settings as in the benchmark simulations without noise in order to obtain results under each observational noise level.

The results of introducing different levels of observational noise in the three scenarios are shown in Figure~\ref{fig_sync_noise1}--\ref{fig_sync_noise3}. In Scenario 1 (Figure~\ref{fig_sync_noise1}) and Scenario 3 (Figure~\ref{fig_sync_noise3}), SCPCM continued to reveal the causal influence between the simulated variables $X$ and $Z$ even as the observational noise increased. In Scenario 2 (Figure~\ref{fig_sync_noise2}), SCPCM began to fail to detect the causal influence between $X$ and $Z$ once the observational noise reached 75\%. At this noise level, the signal to noise ratio was 1 to 3 and the noise overwhelmed the system to a degree that could not be separated by SCPCM.

\begin{figure}[htbp]
\centering
\includegraphics[width=1\textwidth]{figure/SI_sync_noise1.png}
\caption{Cross mapping results under increasing levels of observational noise in benchmark scenario~1.}
\label{fig_sync_noise1}
\end{figure}

\begin{figure}[htbp]
\centering
\includegraphics[width=1\textwidth]{figure/SI_sync_noise2.png}
\caption{Cross mapping results under increasing levels of observational noise in benchmark scenario~2.}
\label{fig_sync_noise2}
\end{figure}

\begin{figure}[htbp]
\centering
\includegraphics[width=1\textwidth]{figure/SI_sync_noise3.png}
\caption{Cross mapping results under increasing levels of observational noise in benchmark scenario~2.}
\label{fig_sync_noise3}
\end{figure}



% \newpage
% \subsection{Impact of spatial detrending on causal discovery}
% \label{sec_detrending}

% Previous studies applying causal discovery algorithms to spatial cross sectional data typically remove a first order spatial trend from the variables based on their geographic locations \citep{gao2023gccm}. We adopt the same detrending strategy in this study and evaluate its influence on the result of SCPCM. Figure~\ref{fig_detrending2} and Figure~\ref{fig_detrending1} present the SCPCM results for two empirical applications in this study without spatial detrending. Except for the infrastructure accessibility and health outcomes case, in which SCPCM correctly identifies the absence of a direct causal influence between social infrastructure accessibility and HALE, SCPCM fails to recover the true causal relationships in the other example. These results indicate that spatial detrending should be performed prior to applying SCPCM to spatial cross sectional data. The functions implemented in our \texttt{spEDM} package therefore enable detrending by default in order to ensure reliable SCPCM result.

% \newline
% \begin{figure}[htbp]
% \centering
% \includegraphics[width=1\textwidth]{figure/SI_detrending2.png}
% \caption{Cross mapping results for farmland net primary productivity case without spatial detrending.}
% \label{fig_detrending2}
% \end{figure}

% \begin{figure}[htbp]
% \centering
% \includegraphics[width=1\textwidth]{figure/SI_detrending1.png}
% \caption{Cross mapping results for infrastructure accessibility and health outcome case without spatial detrending.}
% \label{fig_detrending1}
% \end{figure}

% \newpage
% \subsection{Sensitivity to state-space reconstruction parameters}
% \label{sec_embedding}



\newpage
\section{Best practices for causal discovery with SCPCM}\label{sec_guideline}

\subsection{Brief steps for conducting SCPCM}
\label{sec_steps_scpcm}

In general, SCPCM requires four major steps to complete the causal discovery process:

1.
\textit{Construct spatial embeddings.}  
Set the embedding dimension $E$ and the lag step $\tau$, and identify spatial neighbors of different orders for each spatial unit. Record the values of all variables included in the analysis, such as the driver variables (for example $x$ and $z$) and one or more conditional variables (for example $y$ or other covariates). Organize these values into spatial embedding vectors and assemble them into an embedding matrix according to their lag order.

2.
\textit{Perform direct cross prediction.}  
For a sequence of library sizes, predict $x$ from $z$ and predict $z$ from $x$ using simplex projection. Compute the resulting prediction skills $\rho_{z \; xmap \; x}$ and $\rho_{x \; xmap \; z}$. These two direct prediction quantities represent the overall predictive skills between $x$ and $z$, including both direct and indirect predictive information.

3.
\textit{Perform conditional cross prediction.}  
Repeat the prediction process by augmenting the driver variables with one or more conditional variables. For example, predict $x$ from $z$ while conditioning on a variable $y$ or an extended set of conditioning variables $\mathit{C_s}$. Compute the conditional prediction skills $\rho_{z \; xmap \; x\,|\,\mathit{C_s}}$ and $\rho_{x \; xmap \; z\,|\,\mathit{C_s}}$. These conditional quantities remove the influence of the conditioning variables and allow the direct causal component to be isolated.

4.
\textit{Evaluate the joint behavior of the four prediction quantities and interpret causation.}  
Plot the four cross prediction curves using prediction skill on the vertical axis and library size on the horizontal axis. Assess convergence, significance and confidence intervals of prediction skill with the largest library size, and relative differences between the direct and conditional predictions. SCPCM evaluates the behavior of the four quantities, namely $\rho_{z \; xmap \; x\,|\,\mathit{C_s}}$, $\rho_{x \; xmap \; z\,|\,\mathit{C_s}}$, $\rho_{z \; xmap \; x}$, and $\rho_{x \; xmap \; z}$, across library sizes to determine whether a direct causal influence exists between $x$ and $z$. A persistent improvement in the conditional prediction skills relative to their direct counterparts, accompanied by increasing library size, indicates the presence of a direct causal influence that cannot be explained by the conditioning variables. Conversely, if the conditional prediction deteriorates or fails to increase with library size, SCPCM concludes that the apparent association between $x$ and $z$ is indirect or fully attributable to the conditioning set $\mathit{C_s}$. 

\newpage
\subsection{Practical implementation of spatial causal discovery using the spEDM package}
\label{sec_run_scpcm}

Figure~\ref{fig_best_practices} provides an overview of the recommended procedure for conducting SCPCM based causal discovery using the \texttt{spEDM} package. We developed this user friendly R package to support flexible usage of SCPCM and to integrate it smoothly with the spatial data science ecosystem of R \citep{baseR,rpkg_sf2,geocompr}. The package allows users to load a variety of spatial data through existing tools such as \texttt{readr}, \texttt{readxl}, \texttt{sf} and \texttt{terra}. After data import, the function \texttt{spEDM::fnn()} applies the false nearest neighbours method to identify the minimum embedding dimension $E$. The function \texttt{spEDM::simplex()} is then used to perform univariate self-prediction, where each variable is predicted from its own reconstructed state space, in order to determine the optimal combination of parameters $E$, $\tau$ and $k$. Using this parameter set, users can select the relevant variables and execute \texttt{spEDM::scpcm()} to perform causal discovery with SCPCM on spatial cross sectional data. The output of \texttt{scpcm()} can then be visualized using the built in S3 \texttt{plot} method provided by \texttt{spEDM}, or users may alternatively employ base R or \texttt{ggplot2} to plot the cross prediction skills against library sizes to assess the presence or absence of direct causal influences.

% \vspace{1em}
\begin{figure}[htbp]
\centering
\includegraphics[width=1\textwidth]{figure/SI_scpcm_procedure.png}
\caption{Procedure for applying SCPCM with the \texttt{spEDM} package}
\label{fig_best_practices}
\end{figure}

\newpage
\section*{Supplementary References}
\bibliographystyle{elsarticle-harv}
\bibliography{references.bib} 

\end{document}